% ─────────────────────────────────────────────────────────────────────────────
%  Handout – Erprobung NwT-Lernmodule | Lehrkräfte
%  Veranstaltung: Montag, 07. Dezember 2026 | 15:30 Uhr | DLL am KIT
% ─────────────────────────────────────────────────────────────────────────────
\documentclass[a4paper,10pt]{article}
\usepackage[utf8]{inputenc}
\usepackage[ngerman]{babel}
\usepackage[margin=2cm, top=1.8cm, bottom=2cm]{geometry}
\usepackage{helvet}
\renewcommand{\familydefault}{\sfdefault}
\usepackage{tabularx}
\usepackage{booktabs}
\usepackage{xcolor}
\usepackage{tcolorbox}
\usepackage{enumitem}
\usepackage{graphicx}
\usepackage{fontawesome5}    % optional – nur wenn verfügbar; sonst Zeilen mit \fa... löschen
\usepackage{hyperref}
\hypersetup{
  colorlinks=true,
  linkcolor=primary,
  urlcolor=accent,
  pdfborder={0 0 0}
}

\definecolor{primary}{RGB}{31, 50, 70}
\definecolor{accent}{RGB}{0, 120, 110}
\definecolor{accentlight}{RGB}{220, 245, 242}
\definecolor{warn}{RGB}{180, 60, 40}
\definecolor{lightgray}{RGB}{247, 248, 250}
\definecolor{midgray}{RGB}{140, 150, 160}

\pagestyle{empty}
\setlength{\parindent}{0pt}

\newcommand{\linefill}{\rule{\linewidth}{0.4pt}}
\newcommand{\sectionrule}{%
  \vspace{2pt}{\color{accent}\rule{\linewidth}{1.5pt}}\vspace{4pt}}

% ─────────────────────────────────────────────────────────────────────────────
\begin{document}

% ── KOPFZEILE ─────────────────────────────────────────────────────────────────
{\LARGE \bfseries \color{primary} Erprobung digitaler NwT-Lernmodule}\\[3pt]
{\large \color{accent} Einladung zur Pilotveranstaltung – Rückmeldung aus der Praxis}\\[6pt]
\linefill\\[6pt]

% ── LOGISTIK-BOX ──────────────────────────────────────────────────────────────
\begin{tcolorbox}[
  colback=accentlight,
  colframe=accent,
  arc=2mm,
  boxrule=1pt,
  top=6pt, bottom=6pt, left=10pt, right=10pt,
  title={\bfseries\color{white} Termin \& Ort},
  coltitle=white,
  attach boxed title to top left={yshift=-2mm, xshift=4mm},
  boxed title style={colback=accent, arc=1mm}
]
\begin{tabularx}{\linewidth}{@{}lX@{}}
  \textbf{Datum}   & Montag, 07.\ Dezember 2026 \\[2pt]
  \textbf{Beginn}  & 15:30 Uhr (Ende gegen 17:30 Uhr) \\[2pt]
  \textbf{Ort}     & \textbf{Digital Learning Lab (DLL)} am KIT \\
                   & 1.\ Stock, InformatiKOM (Geb.\ 50.19, Raum 108 / Seminarraum 5) \\
                   & Adenauerring 12, 76131 Karlsruhe \\[2pt]
  \textbf{Für uns} & Snacks \& Getränke werden gestellt – vielen Dank für Ihren Weg!
\end{tabularx}
\end{tcolorbox}

\vspace{10pt}

% ── HINTERGRUND ───────────────────────────────────────────────────────────────
{\large \bfseries \color{primary} Warum diese Veranstaltung?}
\sectionrule

Für NwT-Projekte gibt es Unterrichtseinheiten und Materialien –
doch viele davon sind entweder \textbf{zu oberflächlich} für den konkreten Einsatz
oder nur für eingeloggte Lehrkräfte zugänglich.
Was fehlt, sind \textbf{praxisnahe, digital aufbereitete Lernbegleiter},
die Schülerinnen und Schüler Schritt für Schritt durch ein Projekt führen –
von der Theorie bis zum funktionierenden Ergebnis.

\vspace{6pt}

Genau das sollen die hier vorgestellten \textbf{Lernmodule} leisten:
interaktive Webseiten, die Erklärungen, Simulationen, Codebeispiele
und Aufgaben in einem didaktisch strukturierten Format vereinen.

\vspace{10pt}

% ── DIE MODULE ────────────────────────────────────────────────────────────────
{\large \bfseries \color{primary} Die zwei Module}
\sectionrule

\begin{tcolorbox}[
  colback=lightgray, colframe=midgray, arc=2mm,
  boxrule=0.6pt, top=5pt, bottom=5pt, left=8pt, right=8pt
]
\begin{tabularx}{\linewidth}{@{}lX@{}}
  \textbf{Modul 1} & \textbf{Fotometer} \\
                   & Arduino-basiertes Fotometer: Schaltungsaufbau, Messung von Lichtintensität,
                     Python-Auswertung. \\[6pt]
  \textbf{Modul 2} & \textbf{Wetterstation-App} \\
                   & Netzwerkbasiertes Projekt: Sensoren, Datenübertragung über WLAN,
                     Visualisierung auf einem lokalen Webserver. \\
\end{tabularx}
\end{tcolorbox}

\vspace{4pt}
{\small \color{midgray}
Beide Module sind eigenständig und können unabhängig voneinander bearbeitet werden.}

\vspace{10pt}

% ── ABLAUF ────────────────────────────────────────────────────────────────────
{\large \bfseries \color{primary} Ablauf des Nachmittags}
\sectionrule

\begin{tabularx}{\linewidth}{@{}>{\bfseries}p{2.6cm}X@{}}
  \toprule
  15:30 – 15:45 & Begrüßung, kurze Vorstellung der Lernmodule \\
  15:45 – 17:00 & \textbf{Freies Erkunden:} Wählen Sie ein Modul (oder beide) und arbeiten Sie sich in Ihrem eigenen Tempo vor. Hardware und Laptops stehen bereit. \\
  17:00 – 17:20 & Ausfüllen des Evaluationsfragebogens \\
  17:20 – 17:30 & Gemeinsame Diskussion, Austausch, Fragen \\
  \bottomrule
\end{tabularx}

\vspace{10pt}

% ── IHRE ROLLE ────────────────────────────────────────────────────────────────
{\large \bfseries \color{primary} Ihre Rolle heute}
\sectionrule

Sie kommen mit einem einzigartigen Blickwinkel: Sie kennen NwT-Projekte aus dem Schulalltag.
Ihre Rückmeldung ist deshalb besonders wertvoll – nicht als Lernende,
sondern als \textbf{Expertinnen und Experten für das, was im Unterricht wirklich funktioniert}.

\vspace{6pt}

\begin{itemize}[leftmargin=1.4em, itemsep=2pt]
  \item Probieren Sie die Module aus Lernenden- \emph{und} Lehrendenperspektive aus
  \item Bewerten Sie: Würden Sie ein solches Modul im eigenen Unterricht einsetzen?
  \item Notieren Sie konkrete Verbesserungsideen – sie fließen direkt in die Weiterentwicklung ein
\end{itemize}

\vspace{10pt}

% ── FRAGEBOGEN-HINWEIS ────────────────────────────────────────────────────────
\begin{tcolorbox}[
  colback=lightgray, colframe=midgray, arc=2mm,
  boxrule=0.6pt, top=5pt, bottom=5pt, left=8pt, right=8pt
]
{\bfseries Evaluation:} Am Ende erhalten Sie einen kurzen Fragebogen (ca.\ 15 Min.).
Er ist anonym und fragt nach Verständlichkeit, Motivation, digitalen Elementen
und der Relevanz der Inhalte für Ihren Unterricht.
Ihre ehrliche Rückmeldung hilft dabei, die Module für den Einsatz im
Lehramtsstudium und im Schulunterricht weiterzuentwickeln.
\end{tcolorbox}

\vspace{8pt}

% ── FUSSZEILE ─────────────────────────────────────────────────────────────────
\linefill\\[4pt]
{\small \color{midgray}
Piloterprobung NwT-Lernmodule $\cdot$ Dezember 2026 $\cdot$
Digital Learning Lab (DLL), KIT Karlsruhe
}

\end{document}
